\documentclass[12pt]{article}
\usepackage[utf8]{inputenc}
\usepackage[T1]{fontenc}
\usepackage[margin=1in]{geometry}
\usepackage{setspace}
\usepackage{xcolor}
\usepackage[colorlinks=true, urlcolor=blue!60!black, linkcolor=blue!60!black]{hyperref}
\setlength{\parindent}{0em}
\setlength{\parskip}{0.6em}
\renewcommand{\baselinestretch}{1.0}
\begin{document}
\begin{center}
{\LARGE \textit{Gonzales v. Raich}, 545 U.S. 1 (2005)\par}
\vspace{0.5em}
{\large Case Brief --- Two-Page Edition\par}
\end{center}

\section{Memory Jogger}
Congress may ban homegrown, state-legal medical marijuana because local production of a fungible commodity can undercut a comprehensive interstate drug regime (\href{https://www.courtlistener.com/opinion/799995/gonzales-v-raich/}{CourtListener}).

\section{Detailed Case Facts}
California's Compassionate Use Act of 1996 shielded seriously ill patients who used marijuana on a physician's recommendation. Angel Raich and Diane Monson relied on cannabis as their only effective treatment; Monson grew her own, and Raich received it free from caregivers. The Controlled Substances Act (CSA) makes marijuana a Schedule I substance. In 2002, DEA agents seized and destroyed Monson's six plants, though county officials found her use lawful under state law.

\section{Procedural History}
The district court denied a preliminary injunction. The Ninth Circuit reversed, treating intrastate noncommercial medical use as beyond the commerce power under \textit{Lopez} and \textit{Morrison}. The Supreme Court vacated and remanded.

\section{Judicial Votes}
6--3. Stevens wrote for Kennedy, Souter, Ginsburg, and Breyer; Scalia concurred in the judgment. O'Connor dissented, joined by Rehnquist and Thomas (except Part III); Thomas also dissented separately.

\section{Holding}
The CSA validly applies to intrastate, noncommercial medical marijuana. Congress may regulate a local class of activity when it rationally concludes that exempting it would undercut a comprehensive scheme regulating an interstate market.

\section{Analysis of Opinions}
\textbf{Majority:} As in \textit{Wickard v. Filburn}, home production of a fungible commodity affects the national market, and exempting it would leave a ``gaping hole'' in the CSA. \textit{Lopez} and \textit{Morrison} involved single-subject statutes regulating noneconomic conduct. \textbf{Scalia:} the Necessary and Proper Clause lets Congress reach even noneconomic local activity to make interstate regulation effective. \textbf{O'Connor:} the comprehensive-scheme theory makes \textit{Lopez} a drafting guide and the definition of ``economic'' is ``breathtaking.'' \textbf{Thomas:} on this logic Congress may ``regulate quilting bees, clothes drives, and potluck suppers throughout the 50 States.''

\section{Examples: Future Applications}
\textbf{Same-Side:} a comprehensive firearms statute reaches machine guns assembled at home. \textbf{Opposite-Side:} a single-subject ban on backyard composting near schools fails like \textit{Lopez}. \textbf{Fence-Sitter:} the Endangered Species Act protects a one-county cave beetle that is no fungible commodity.

\section{Critique}
Originalists such as Randy Barnett say \textit{Raich} abandoned the founding meaning of ``commerce,'' and \textit{NFIB v. Sebelius} later cabined it. Others note the irony of liberals defending the War on Drugs and the muddled rule that local regulation survives inside larger schemes.

\section{Key Quotations}
\begin{quote}
``We need not determine whether respondents' activities, taken in the aggregate, substantially affect interstate commerce in fact, but only whether a `rational basis' exists for so concluding.''
\end{quote}
\noindent --- Justice Stevens, for the Court, 545 U.S. 1.

\end{document}
